% PROVENANCE (section-level)
%   produced-by : 5_BBbkg_weights_optuna_minuit.py, bbbar_reweighting.py
%   commit      : ad2f2ec
\section{Generic-$B\bar{B}$ Background Modelling and Validation}
\label{sec:bbbar}

Generic $B\bar{B}$ events in which a true $D$ and a true lepton are
reconstructed --- either from different $B$ mesons
(\texttt{bkg\_combinatorial}) or from a single hadronic $B$ decay supplying a
secondary lepton (\texttt{bkg\_hadronicB\_secondaryL}) --- are the dominant
background and are mismodelled in MC16rd.
They are corrected by a data-driven reweighting of the underlying $B$-decay
composition.

\subsection{Hierarchical decay classification}
\label{sec:bbbar:hierarchy}
For each of the two categories, the generated decay of the $B$ that supplied
the $D$ (and optionally of the $B$ that supplied the lepton) is identified by
the product of its daughter PDG codes, and assigned to one family:
\begin{enumerate}
  \item \texttt{BBbar\_semileptonic} --- any generated daughter is a lepton;
  \item \texttt{BBbar\_measured\_hadronic} --- the decay-mode identifier
        appears in the list of modes with a measured branching fraction,
        in which case a PDG branching-fraction correction is also applied;
  \item \texttt{BBbar\_unmeasured:$n$-body} --- otherwise, grouped by the
        generated daughter multiplicity, capped at $5+$.
\end{enumerate}
The hierarchy is applied in this order, so the families are mutually
exclusive and, because the last is a catch-all over multiplicity, exhaustive.
Its purpose is to separate what is externally constrained (measured branching
fractions) from what is not (the unmeasured hadronic modes, which carry the
bulk of the modelling freedom), and to group the unmeasured modes by a
variable --- multiplicity --- that drives the ROE observables without
requiring a per-mode measurement.

\subsection{Family normalisation versus internal composition}
\label{sec:bbbar:normalisation}
One weight is fitted per family; the composition \emph{within} a family is
left at its generated values.
This is deliberate: floating individual decay modes without external
constraints is not supported by the available statistics.
An optional many-to-one mode replacement merges configured non-resonant
3- and 4-body modes into resonant 2-body states before tuning, preserving the
total yield of the affected modes through an $(N_{\mathrm{old}} +
N_{\mathrm{new}})/N_{\mathrm{new}}$ rescaling of the surviving mode.
It is disabled in the results shown here (\texttt{replaceFalse}).

The reweighting does \emph{not} preserve the inclusive $B$-decay rate, and is
not intended to.
The generic-$B\bar{B}$ normalisation differs between data and MC in the
tuning region; requiring the inclusive rate to remain fixed would force that
discrepancy entirely into the composition, which is the quantity being
measured.
The total rate is therefore allowed to move, and the requirement is instead
that the resulting change be validated in a region independent of the one it
was tuned in (Sec.~\ref{sec:bbbar:validation}).

\subsection{Tuning procedure and region}
\label{sec:bbbar:tuning}
Five family weights are fitted by minimising a binned Poisson deviance
between data and weighted MC, using Optuna (TPE sampler, \texttt{MedianPruner})
for global exploration followed by iminuit MIGRAD/HESSE/MINOS.
\texttt{BBbar\_semileptonic} is held at unity and the fake-$D$ normalisation
is held at its externally determined value, so they are not free parameters.
The set of fit bins is derived once from unweighted MC support and held
fixed, so that the likelihood support cannot depend on the parameters being
fitted.

Three objectives are implemented in \texttt{5\_BBbkg\_weights\_optuna\_minuit.py}
and have all been tried.
The nominal objective is \texttt{kinematic-2d}: a two-dimensional deviance in
$\mmiss \times \pdl$ ($20 \times 20$ bins over
$[-4,10]\,\mathrm{GeV}^{2}/c^{4} \times [0.2,4.0]\,\mathrm{GeV}/c$) alone
(\texttt{ANALYSIS\_DECISION}).
\texttt{kinematic-2d-plus-roe} adds a shape-only one-dimensional deviance in
the ROE track multiplicity, weighted by a coefficient \texttt{roe\_strength}
(``$\alpha$''), with the final bin merging all multiplicities $\geq 14$;
\texttt{missm2-roe-2d} instead replaces the $\mmiss \times \pdl$ deviance with
a two-dimensional missing-mass/ROE-multiplicity deviance.
Adding the ROE-multiplicity term (at $\alpha = 0.5$ or $\alpha = 1$; $\alpha$
itself has not been tuned, only tried at these two values) does not improve
the agreement on the $\mmiss \times \pdl$ plane, and \texttt{missm2-roe-2d}
performs worse still, both in total deviance and in how pinned the fitted
parameters are (Sec.~\ref{sec:bbbar:results}).
One interpretation is that the ROE tracks and the two kinematic variables
constrain different aspects of the generic-$B\bar{B}$ composition, so that
adding the ROE term does not simply sharpen the same likelihood but pulls it
towards a different, not obviously more correct, optimum;
\texttt{kinematic-2d} is used as the nominal objective on the grounds that it
gives the best agreement on the $\mmiss \times \pdl$ plane specifically, since
that is the plane the final fit is performed in
(Sec.~\ref{sec:fit:observables}).
\AnalysisTBD{the interpretation above (different physics constrained by the
ROE tracks versus the two kinematic variables) is a hypothesis, not yet
demonstrated}

The tuning region is
\begin{equation}
  1.855 < m(K\pi\pi) < 1.885,\quad
  \texttt{fakeD\_prob} < 0.1,\quad
  \texttt{sig\_prob} < 0.1 ,
  \label{eq:bbbar-region}
\end{equation}
together with the offline requirements of Eq.~\eqref{eq:offline-cut} except
$p_{\ell} < 4$ (Sec.~\ref{sec:recon:offline}).
Note that Eq.~\eqref{eq:bbbar-region} is a \emph{signal-region} $D$-mass
selection with the signal and fake-$D$ classifier outputs suppressed; it is
not the combinatorial-enriched region of \texttt{lgb\_comb}, and it places no
requirement on \texttt{combinatorial\_prob} or \texttt{continuum\_prob}
(Appendix~\ref{app:selections:naming}).

\subsection{Run dependence}
\label{sec:bbbar:rundep}
Weights are derived separately for each run period and channel.
For the combined \texttt{run1+run2} fit the differing fake-$D$ normalisations
($0.87$ for Run~1, $1.0$ for Run~2) are applied as an event-level factor
rather than as a common category weight, so that the run dependence survives
the combination.

\subsection{Current status of the tuned weights}
\label{sec:bbbar:results}

The stored fits were regenerated for the combined \texttt{run1+run2} sample,
and MINOS was enabled for every scenario, specifically to address the
parameter-at-limit problem below.
Six results are stored, all with mode replacement disabled
(\texttt{replaceFalse}): \texttt{kinematic-2d} for $e$ and $\mu$;
\texttt{kinematic-2d-plus-roe} at $\alpha=0.5$ for $e$ and $\mu$ and at
$\alpha=1$ for $e$; and \texttt{missm2-roe-2d} for $e$.
The nominal result is \texttt{kinematic-2d} (Sec.~\ref{sec:bbbar:tuning}).

% PROVENANCE
%   status      : PRELIMINARY_RESULT
%   produced-by : BBbkg_weights/*.json (fit metadata read directly)
%   commit      : ad2f2ec
%   sample      : MC16rd + Proc16/Prompt16, Run 1 + Run 2, e and mu
%   selection   : tuning region of Eq. (bbbar-region)
%   corrections : PID, generic-MC luminosity 0.25, fixed fake-D and
%                 semileptonic weights
\PreliminaryResult{nominal (\texttt{kinematic-2d}, run1+run2):
$e$ channel --- measured hadronic $2.264^{+0.057}_{-0.057}$, unmeasured
2-body $0.0010^{+0.158}_{-0.000}$ (at its lower bound of $0.001$), 3-body
$0.765^{+0.114}_{-0.118}$, 4-body $0.0010^{+0.031}_{-0.000}$ (at its lower
bound), 5+-body $0.027^{+0.046}_{-0.026}$; total deviance $419.78$.
$\mu$ channel --- measured hadronic $1.925^{+0.111}_{-0.112}$, unmeasured
2-body $0.0010^{+0.088}_{-0.000}$ (at its lower bound), 3-body
$0.583^{+0.202}_{-0.204}$, 4-body $0.0010^{+0.112}_{-0.000}$ (at its lower
bound), 5+-body $0.046^{+0.065}_{-0.045}$; total deviance $430.07$.
Uncertainties quoted are MINOS; parameter bounds are
$[0.5,5]$ for the measured-hadronic weight and $[0.001,5]$,
$[0.01,1]$, $[0.001,1]$, $[0.001,0.5]$ for the 2-, 3-, 4- and 5+-body
unmeasured weights respectively (\texttt{PARAMETER\_SPECS} in
\texttt{5\_BBbkg\_weights\_optuna\_minuit.py}).  Source:
\texttt{BBbkg\_weights/best\_bbbar\_weights\_kinematic-2d\_run1+run2\_\{e,mu\}%
\_replaceFalse\_minuit.json}}

\InProgress{turning on MINOS improves the situation but does not resolve it.
In both channels of the nominal fit, and in every one of the six stored
fits, the unmeasured 2-body and 4-body weights remain pinned at their lower
bound (\texttt{has\_parameters\_at\_limit = true} throughout, exactly as
before MINOS was enabled); the parameter bounds themselves are unchanged.
For the three parameters that are \emph{not} pinned in a given fit ---
typically the measured-hadronic, 3-body and 5+-body weights --- MINOS does
give a genuine, asymmetric confidence interval that HESSE could not, and
those three could in principle already supply a propagatable width.
For the two parameters that remain pinned, however, the MINOS interval is not
a resolution of the boundary problem: its lower bound is, by construction,
at or immediately next to the parameter's hard limit (the reported lower
errors above are $\mathcal{O}(10^{-7})$ or smaller), so it reflects the
imposed bound rather than a profiled uncertainty, and the same
three-way question of Sec.~\ref{sec:bbbar:degeneracy} --- too-tight bounds,
an unresolved degeneracy, or a genuinely non-physical preferred weight ---
remains open for exactly those two families.  The \texttt{missm2-roe-2d}
result pins a third parameter (3-body, in addition to 2-body) and pushes
5+-body to its \emph{upper} bound instead, which is consistent with the
qualitative statement above that this objective performs worse, not better}

\subsection{Constraining power of the tuning region}
\label{sec:bbbar:degeneracy}

The tuned fits do not reach an interior minimum.
In all six results currently stored in \texttt{BBbkg\_weights/} the minimum
sits on at least one parameter bound, and the pinned parameters are always
among the \emph{unmeasured} $n$-body weights (2-body and 4-body in five of
the six fits; \texttt{missm2-roe-2d} additionally pins 3-body and pushes
5+-body to its \emph{upper} bound) and never the measured-hadronic weight.
Table~\ref{tab:bbbar-correlations} shows the accompanying correlation
structure.
The pairwise correlation between the 3-body and 5+-body unmeasured weights is
large in every fit that uses the $\mmiss\times\pdl$-based objectives, $0.75$
to $0.78$; \texttt{missm2-roe-2d} is the exception, with unmeasured pairwise
correlations below $0.001$ but instead a moderate correlation of the
measured-hadronic weight with 3-body ($0.30$).
In the other five fits the measured-hadronic weight also correlates with
3-body and 5+-body at the $0.7$--$0.8$ level, larger than in the fits stored
previously, though it still never pins.

% PROVENANCE
%   status      : PRELIMINARY_RESULT (fit metadata, not a physics measurement)
%   produced-by : BBbkg_weights/*.json (fit metadata read directly)
%   commit      : ad2f2ec
%   sample      : MC16rd + Proc16/Prompt16, Run 1 + Run 2, e and mu
%   selection   : tuning region of Eq. (bbbar-region)
%   corrections : PID, generic-MC luminosity 0.25, fixed fake-D and
%                 semileptonic weights
\begin{table}[htbp]
  \centering
  \caption{Correlation structure of the stored generic-$B\bar{B}$ weight fits,
  all \texttt{run1+run2}, mode replacement disabled.
  ``unmeas.\ pair'' is the largest absolute correlation between two unmeasured
  $n$-body weights; ``meas.'' is the largest absolute correlation of the
  measured-hadronic weight with any other parameter.
  Every fit reported \texttt{has\_parameters\_at\_limit = true}.}
  \label{tab:bbbar-correlations}
  \scriptsize
  \setlength{\tabcolsep}{3pt}
  \begin{tabular}{>{\raggedright\arraybackslash}p{0.30\textwidth} >{\centering\arraybackslash}p{0.10\textwidth} >{\centering\arraybackslash}p{0.10\textwidth} >{\raggedright\arraybackslash}p{0.32\textwidth}}
    \hline
    Fit & $\max|\rho|$\newline unmeas. & $\max|\rho|$\newline meas. & Pinned \\
    \hline
    \texttt{kinematic-2d $e$} (nominal) & 0.77 & 0.71 & 2-body, 4-body \\
    \texttt{kinematic-2d $\mu$} (nominal) & 0.75 & 0.79 & 2-body, 4-body \\
    \texttt{kin-2d-plus-roe $\alpha{=}0.5$ $e$} & 0.78 & 0.75 & 2-body, 4-body \\
    \texttt{kin-2d-plus-roe $\alpha{=}0.5$ $\mu$} & 0.76 & 0.81 & 2-body, 4-body \\
    \texttt{kin-2d-plus-roe $\alpha{=}1$ $e$} & 0.78 & 0.77 & 2-body, 4-body \\
    \texttt{missm2-roe-2d $e$} & $<0.001$ & 0.30 & 2-body, 3-body, 5+-body (upper) \\
    \hline
  \end{tabular}
\end{table}

Three explanations are consistent with this pattern, and they call for
different responses.
\begin{enumerate}
  \item The parameter bounds are simply too tight, in which case relaxing
        them gives an interior minimum and a usable covariance.
  \item The tuning region cannot separate the unmeasured families, in which
        case the likelihood is flat along at least one direction and relaxing
        the bounds only lets the minimum slide further before pinning again.
        The remedy is to merge the families the region cannot separate, or to
        add an observable that separates them.
  \item The constrained optimum is genuine: the data prefer a weight at or
        below zero for one or more unmeasured families, so the optimum lies
        outside the allowed range.  Relaxing the bounds does not help and
        merging families does not either; the question becomes why a decay
        family is preferred at a non-physical rate, which points at the
        family definition or at the simulation rather than at the fit.
\end{enumerate}

Cases 2 and 3 are not distinguished by the deviance alone: if the optimum lies
below every relaxed lower bound, the deviance improvement between successive
relaxations tends to zero as that bound tends to zero, which mimics a flat
direction.  They are separated by profiling a pinned weight inward while
re-minimising the others --- flat for a genuine degeneracy, rising for a
constrained optimum.  A one-parameter scan does not suffice, because moving one
weight along a degenerate direction while holding the others fixed also raises
the objective.
That the measured-hadronic weight --- the one family with an external
constraint on its composition --- never pins, while the unconstrained families
repeatedly do, disfavours the first explanation and is consistent with either
the second or the third.
The strong correlations of Table~\ref{tab:bbbar-correlations} point towards
the second, but they do not exclude the third, and the distinction matters
because the two call for opposite responses.
Turning on MINOS (Sec.~\ref{sec:bbbar:results}) does not by itself distinguish
the three either: a MINOS interval that touches the bound on a pinned
parameter is exactly what both the second and the third explanation predict,
and does not test the first.
\InProgress{the three explanations are distinguished by re-running the
minimisation over a sequence of bound relaxations and profiling each pinned
weight inward while re-minimising the others.
\texttt{scripts/scan\_bbbar\_bounds.py} in the analysis repository does this
and is now included at the commit pinned by this note, but it has still never
been run on real ntuples.  Everything above is therefore an inference from
the metadata of six stored fits.  No conclusion about merging decay families
should be drawn until the scan has been run and its output cross-checked}

\subsection{Validation}
\label{sec:bbbar:validation}
\AnalysisTBD{an independent validation region for the tuned weights has not
been established.  The wrong-charge sample was the candidate, but its physics
composition differs from that of the tuning region --- different hadronic
versus semileptonic fractions, and a different hadronic multiplicity spectrum
(see \texttt{Notebooks/Plotting/4\_Plot\_WC\_*.ipynb}) --- so agreement or
disagreement there does not cleanly test the weights.  A validation region
matched in composition to the tuning region is needed}

\AnalysisTBD{closure of the tuned weights in the bins actually used by the
final fit (Sec.~\ref{sec:fit}), which are not the same bins as the tuning
histogram}
